\chapter{Appendix: Complex Numbers, Linear Algebra, and Calculus for DSP}

The body of this book leans on three branches of mathematics so often, and so casually, that it
is worth gathering the essential pieces into one place where you can find them quickly. This
appendix is that place. It is not a substitute for the companion volumes -- \emph{Calculus From
the Ground Up}, \emph{Linear Algebra From the Ground Up}, and \emph{Probability From the Ground
Up} -- each of which builds these ideas slowly and from scratch. The goal here is narrower and
more practical: to give you, in one sitting, exactly the geometric pictures and exactly the
formulas that signal processing actually uses, with enough intuition that the symbols stop being
intimidating and start being \emph{useful}.

We will work in the same spirit as the rest of the book: geometry first, formula second. A
complex number is a point in a plane and, more importantly, a little machine for rotating. A
vector is an arrow; a matrix is a way of moving arrows around; an inner product measures how much
two arrows agree. An integral is a limit of sums that adds up a signal continuously. A
probability distribution is a way of describing the cloud of values that noise can take. None of
this is hard once you see it, and all of it is the connective tissue of DSP.

\section{Why this toolkit, and how the pieces fit}
\label{ch42:overview}

Before the details, here is the map. Three threads run through every chapter of this book, and
each one corresponds to a section below.

\begin{itemize}
  \item \textbf{Complex numbers} (Section~\ref{ch42:complex}) are how we describe \emph{rotation
  at a steady rate}, which is to say frequency. The single object $e^{\jj\omega t}$ -- a point
  marching around a circle -- replaces the awkward pair $(\cos,\sin)$ and turns trigonometry into
  arithmetic. Without it, the Fourier transform is a mess of sum-to-product identities; with it,
  the whole theory becomes one clean idea.
  \item \textbf{Linear algebra} (Section~\ref{ch42:linalg}) is how we treat a signal as a single
  geometric object -- a vector -- and how we measure, project, and decompose it. The Fourier
  transform is nothing but a change of basis; the DFT is literally a matrix; filtering is a
  linear operator. Inner products, orthonormal bases, and eigenvectors are the load-bearing
  concepts.
  \item \textbf{Calculus} (Section~\ref{ch42:calculus}) is how we add up a continuous signal
  (integration, which powers the continuous-time Fourier and Laplace transforms) and how we
  handle the idealized spike called the delta function. We need only a little, but we need it
  precisely.
  \item \textbf{Probability} (Section~\ref{ch42:probability}) is how we describe noise: the
  random part of every real measurement. Just enough -- mean, variance, the Gaussian, white
  noise -- to read the noise chapters and to reason about an SNR.
\end{itemize}

\begin{focusbox}
Read this appendix the way you would read a phrasebook before a trip: skim it once so you know
what is here, then come back to the relevant page when the main text uses something you want
nailed down. Every section ends by pointing to the companion book where the idea is built
properly.
\end{focusbox}

\section{Complex numbers and Euler's formula}
\label{ch42:complex}

\subsection{A number that is also a point}

A real number lives on a line. A \emph{complex number} lives in a plane. We write
\[
  z = a + \jj\,b,
\]
where $a$ and $b$ are ordinary real numbers and $\jj$ is a new symbol with the single defining
property
\[
  \jj^2 = -1.
\]
(Engineers write $\jj$ rather than $i$ because $i$ is already taken by electric current; this
book uses $\jj$ throughout.) We call $a = \operatorname{Re}(z)$ the \emph{real part} and $b =
\operatorname{Im}(z)$ the \emph{imaginary part}. The point is to stop thinking of $\jj$ as
mysterious and start thinking of $z$ as the point $(a,b)$ in a two-dimensional plane -- the
\emph{complex plane}, with the real part along the horizontal axis and the imaginary part along
the vertical.

\begin{intuitionbox}
The word ``imaginary'' is a historical accident and an active obstacle to understanding. There is
nothing imaginary about the vertical axis. A complex number is just an ordered pair of reals with
a special multiplication rule -- and that rule, as we are about to see, encodes rotation. Read
$\jj$ as ``the quarter turn'' and the fog lifts.
\end{intuitionbox}

\subsection{Arithmetic: addition slides, multiplication rotates}

Addition is componentwise, exactly like adding vectors:
\[
  (a + \jj b) + (c + \jj d) = (a+c) + \jj(b+d).
\]
Geometrically, adding $z_2$ to $z_1$ slides $z_1$ over by the arrow $z_2$. Nothing surprising.

Multiplication is where the magic is. Using $\jj^2 = -1$ and ordinary distribution,
\[
  (a + \jj b)(c + \jj d) = (ac - bd) + \jj(ad + bc).
\]
This looks like a mess of four products, but its \emph{geometric} meaning is stunningly simple,
and to see it we change coordinates.

\subsection{Polar form: magnitude and angle}

Every point in the plane can be named two ways: by its rectangular coordinates $(a,b)$, or by how
far it is from the origin and in what direction. The distance is the \emph{magnitude} (or
modulus),
\[
  r = \abs{z} = \sqrt{a^2 + b^2},
\]
and the direction is the \emph{argument} (or phase angle),
\[
  \theta = \arg(z) = \operatorname{atan2}(b, a),
\]
measured counterclockwise from the positive real axis. (Use \texttt{atan2}, not plain
\texttt{atan}, so that the quadrant is correct.) Then
\[
  a = r\cos\theta, \qquad b = r\sin\theta, \qquad
  z = r(\cos\theta + \jj\sin\theta).
\]

\begin{figure}[H]
\centering
\begin{tikzpicture}
\begin{axis}[
    width=8.5cm, height=8.5cm,
    axis equal image,
    xlabel={$\operatorname{Re}$}, ylabel={$\operatorname{Im}$},
    xmin=-1.6, xmax=1.6, ymin=-1.6, ymax=1.6,
    axis lines=middle, xtick={-1,1}, ytick={-1,1},
    samples=120, domain=0:360,
]
% unit circle
\addplot[gray, thick] ({cos(x)}, {sin(x)});
% the number z = 0.95 + j0.6 roughly r=1.12 theta~32deg, draw on circle scaled
\addplot[sourceblue, ultra thick, -{Latex}] coordinates {(0,0) (1.15,0.72)};
% projections
\addplot[dashed, gray] coordinates {(1.15,0) (1.15,0.72)};
\addplot[dashed, gray] coordinates {(0,0.72) (1.15,0.72)};
% angle arc
\addplot[specred, thick, domain=0:32, samples=30] ({0.4*cos(x)}, {0.4*sin(x)});
\node[sourceblue] at (axis cs:0.75,0.62) {$z$};
\node[specred] at (axis cs:0.55,0.13) {$\theta$};
\node[gray] at (axis cs:1.15,-0.13) {$a$};
\node[gray] at (axis cs:-0.13,0.72) {$b$};
\node[sourceblue] at (axis cs:0.45,0.45) {$r$};
\end{axis}
\end{tikzpicture}
\caption{A complex number $z = a + \jj b$ as a point in the plane. Its magnitude $r$ is the
length of the arrow from the origin; its argument $\theta$ is the angle from the real axis. The
gray circle is the unit circle $\abs{z}=1$.}
\label{ch42:fig-polar}
\end{figure}

Now watch what multiplication does in polar form. If $z_1 = r_1(\cos\theta_1 + \jj\sin\theta_1)$
and $z_2 = r_2(\cos\theta_2 + \jj\sin\theta_2)$, then expanding the product and collecting terms
using the angle-addition identities for sine and cosine gives
\[
  z_1 z_2 = r_1 r_2\big(\cos(\theta_1 + \theta_2) + \jj\sin(\theta_1 + \theta_2)\big).
\]

\begin{happeningbox}
Multiplying two complex numbers \textbf{multiplies their magnitudes and adds their angles}. So
multiplication is ``scale and rotate.'' In particular, multiplying by a number of magnitude $1$
is a \emph{pure rotation} -- it spins the plane without stretching it. And multiplying by $\jj$
itself (magnitude $1$, angle $90^\circ$) rotates everything a quarter turn counterclockwise. That
is why $\jj^2 = -1$: two quarter turns is a half turn, which sends $1$ to $-1$. The defining
equation was geometry all along.
\end{happeningbox}

\subsection{Euler's formula: the heart of the matter}

The polar form begs for a shorthand, and there is one of breathtaking economy. Euler's formula
states
\[
  \boxed{\,e^{\jj\theta} = \cos\theta + \jj\sin\theta.\,}
\]
The left side is the exponential function evaluated at an imaginary argument; the right side is
the point on the unit circle at angle $\theta$. With it, the polar form collapses to
\[
  z = r\,e^{\jj\theta},
\]
and the multiplication rule becomes the ordinary law of exponents:
\[
  z_1 z_2 = r_1 e^{\jj\theta_1}\,r_2 e^{\jj\theta_2} = r_1 r_2\, e^{\jj(\theta_1+\theta_2)}.
\]
``Add the angles'' is literally ``add the exponents.'' This is why DSP is written in complex
exponentials and not in cosines and sines: the algebra is trivial.

\begin{whybox}
Why should $e^{\jj\theta}$ trace a circle? Think about velocity. Differentiate $f(\theta) =
e^{\jj\theta}$ with respect to $\theta$: by the chain rule $f'(\theta) = \jj\,e^{\jj\theta} =
\jj\,f(\theta)$. Multiplying by $\jj$ is a $90^\circ$ rotation, so the velocity is always
exactly perpendicular to the position vector, and (since $\abs{\jj}=1$) always the same length.
A point whose velocity is perpendicular to its radius and constant in size can only be going
around a circle at constant speed. Euler's formula is the statement ``rotation is exponential
growth turned sideways.''
\end{whybox}

Two consequences we use constantly. Setting $\theta = \pi$ gives \emph{Euler's identity},
$e^{\jj\pi} + 1 = 0$, linking the five most famous constants. Far more useful for us, adding and
subtracting $e^{\jj\theta}$ and $e^{-\jj\theta}$ recovers the real trig functions:
\[
  \cos\theta = \frac{e^{\jj\theta} + e^{-\jj\theta}}{2}, \qquad
  \sin\theta = \frac{e^{\jj\theta} - e^{-\jj\theta}}{2\jj}.
\]
These \emph{inverse Euler formulas} are how a real sinusoid splits into a pair of counter-rotating
phasors, one at frequency $+\omega$ and one at $-\omega$ -- the origin of the two-sided spectrum
(Figure~\ref{ch42:fig-counterrot}).

\begin{figure}[H]
\centering
\begin{tikzpicture}
  \draw[gray, thick] (0,0) circle (1.4);
  \draw[->, gray] (-1.8,0) -- (1.8,0) node[right] {$\operatorname{Re}$};
  \draw[->, gray] (0,-1.8) -- (0,1.8) node[above] {$\operatorname{Im}$};
  % two conjugate phasors at +/- 40 deg
  \draw[sourceblue, ->, ultra thick] (0,0) -- (40:1.4);
  \node[sourceblue] at (40:1.7) {$\tfrac12 e^{\jj\theta}$};
  \draw[specred, ->, ultra thick] (0,0) -- (-40:1.4);
  \node[specred] at (-40:1.7) {$\tfrac12 e^{-\jj\theta}$};
  % resultant sum lands on real axis
  \coordinate (s) at ({2*1.4*cos(40)*0.5},0);
  \draw[quotientorange, ->, ultra thick] (0,0) -- (s);
  \node[quotientorange] at ($(s)+(0,-0.3)$) {$\cos\theta$};
  % rotation arrows
  \draw[sourceblue, ->] (55:1.0) arc (55:75:1.0);
  \draw[specred, ->] (-55:1.0) arc (-55:-75:1.0);
\end{tikzpicture}
\caption{A cosine as the sum of two counter-rotating half-amplitude phasors. As $\theta$ grows,
the blue phasor spins counterclockwise (frequency $+\omega$) and the red one clockwise (frequency
$-\omega$); their imaginary parts always cancel, so the sum stays on the real axis and traces
$\cos\theta$. This is why a real tone shows up at both $+\omega$ and $-\omega$ in a spectrum.}
\label{ch42:fig-counterrot}
\end{figure}

\begin{figure}[H]
\centering
\begin{tikzpicture}
\begin{axis}[
    width=12cm, height=5.2cm,
    xlabel={$\theta$ (radians)}, ylabel={value},
    xmin=0, xmax=6.283, ymin=-1.3, ymax=1.3,
    axis lines=left, samples=160, domain=0:6.283,
    xtick={0,1.571,3.142,4.712,6.283},
    xticklabels={$0$,$\pi/2$,$\pi$,$3\pi/2$,$2\pi$},
    legend pos=south west, legend cell align=left,
]
\addplot[sourceblue, thick] {cos(deg(x))};
\addlegendentry{$\operatorname{Re}\,e^{\jj\theta}=\cos\theta$}
\addplot[specred, thick, dashed] {sin(deg(x))};
\addlegendentry{$\operatorname{Im}\,e^{\jj\theta}=\sin\theta$}
\end{axis}
\end{tikzpicture}
\caption{The real and imaginary parts of $e^{\jj\theta}$ as $\theta$ runs once around the circle.
The shadow on the real axis is a cosine; the shadow on the imaginary axis is a sine. A sinusoid
is the shadow of a spinning point.}
\label{ch42:fig-euler-shadow}
\end{figure}

\subsection{Conjugate, magnitude squared, and division}

The \emph{complex conjugate} of $z = a + \jj b = r e^{\jj\theta}$ is
\[
  \bar z = a - \jj b = r e^{-\jj\theta},
\]
its mirror image across the real axis -- same magnitude, opposite angle. Conjugation is the
single most useful operation for extracting real quantities from complex expressions, because
\[
  z\,\bar z = (a+\jj b)(a - \jj b) = a^2 + b^2 = \abs{z}^2.
\]
Multiplying a number by its conjugate kills the imaginary part and leaves the magnitude squared.
We use this everywhere: a power spectrum is $\abs{X(f)}^2 = X(f)\,\overline{X(f)}$, and the energy
of a complex signal sample is $\abs{x}^2 = x\bar x$. Division is then conjugate-times-reciprocal:
\[
  \frac{z_1}{z_2} = \frac{z_1\,\bar z_2}{\abs{z_2}^2}
  = \frac{r_1}{r_2}\,e^{\jj(\theta_1 - \theta_2)}.
\]

\begin{example}[Computing with both forms]
\label{ch42:ex-arith}
Let $z = 1 + \jj$. Rectangular gives $a = b = 1$. Polar: $r = \sqrt{1^2 + 1^2} = \sqrt 2$ and
$\theta = \operatorname{atan2}(1,1) = \pi/4$, so $z = \sqrt 2\,e^{\jj\pi/4}$. Then squaring is
trivial in polar form: $z^2 = (\sqrt 2)^2 e^{\jj\pi/2} = 2\jj$. Check in rectangular: $(1+\jj)^2
= 1 + 2\jj + \jj^2 = 1 + 2\jj - 1 = 2\jj$. The conjugate is $\bar z = 1 - \jj$, and $z\bar z =
(1)^2 + (1)^2 = 2 = \abs{z}^2$.
\end{example}

\begin{example}[A phasor over time]
\label{ch42:ex-phasor}
A signal $x(t) = A e^{\jj(\omega t + \phi)}$ is a phasor of fixed length $A$ whose angle
increases linearly at rate $\omega$ rad/s, starting at phase $\phi$. Its real part $A\cos(\omega
t + \phi)$ is the physical sinusoid you would measure. Sampling at intervals $T_s$ replaces $t$
by $nT_s$, giving $x[n] = A e^{\jj\phi}\,(e^{\jj\omega T_s})^n$: a geometric sequence whose
common ratio sits on the unit circle. That single fact -- sampled sinusoids are powers of a
unit-circle number -- is the seed of the entire $z$-transform.
\end{example}

\begin{connectionbox}
The full story of complex numbers as rotations, of $e^{\jj\theta}$ derived from its power series,
and of the complex plane as a field, lives in the early chapters of \emph{Calculus From the
Ground Up} (exponentials and the number $e$) and in the rotation chapter of \emph{Linear Algebra
From the Ground Up}. Section~\ref{ch42:complex} of this appendix is the working summary; those
books are where it is proven.
\end{connectionbox}

\begin{goalbox}
On the wrist tracker, complex numbers show up the instant you compute a spectrum of the IMU or
EMG stream. When you run an FFT on a window of accelerometer samples to find the dominant
frequency of a tremor or a repetitive gesture, every output bin is a complex number $r
e^{\jj\theta}$: $r$ tells you how much energy sits at that frequency (useful), $\theta$ tells you
the phase (often discarded for gesture features, but essential if you later want to align two
sensors). Reading a magnitude spectrum is just taking $\abs{z}$ of a list of complex numbers.
\end{goalbox}

\section{Vectors, matrices, inner products, bases, and eigenvectors}
\label{ch42:linalg}

\subsection{A signal is a vector}

The central conceptual move of this whole book -- made first in Chapter~3 of the main text -- is
to regard a signal as a single point in a high-dimensional space. A length-$N$ discrete signal
$x[0], x[1], \dots, x[N-1]$ is the vector
\[
  \vec x = \big(x[0], x[1], \dots, x[N-1]\big) \in \C^N.
\]
Each \emph{sample} is one coordinate. A two-second audio clip at $44.1$ kHz is a point in a space
of about $88{,}000$ dimensions. We cannot picture that, but every operation we want -- adding
signals, scaling them, measuring their length, asking whether two are ``perpendicular'' --
behaves exactly as it does for arrows in the plane. That is the whole reason to use the language
of vectors: our two-dimensional intuition transfers, unchanged, to a million dimensions.

\subsection{Inner product: how much two signals agree}

For real vectors the inner product (dot product) is the sum of products of corresponding
coordinates. For complex vectors -- and DSP is full of them -- we must conjugate one factor so
that a vector's inner product with itself comes out real and nonnegative:
\[
  \inner{\vec x}{\vec y} = \sum_{n=0}^{N-1} x[n]\,\overline{y[n]}.
\]
The conjugate is not a nuisance; it is the rule that makes geometry work in $\C^N$. With it,
\[
  \norm{\vec x}^2 = \inner{\vec x}{\vec x} = \sum_{n=0}^{N-1} \abs{x[n]}^2,
\]
the total energy of the signal -- a real, nonnegative number, as a squared length must be. The
norm $\norm{\vec x}$ is the signal's length.

\begin{intuitionbox}
The inner product answers one question (Figure~\ref{ch42:fig-projection}): \emph{how much of
$\vec y$ is present in $\vec x$?} If the
two point the same way it is large and positive; if they are perpendicular it is zero; if
opposite, negative. The Fourier transform asks exactly this question, with $\vec y$ ranging over
pure sinusoids of every frequency: ``how much of \emph{this} frequency is in my signal?'' Each
Fourier coefficient is one inner product.
\end{intuitionbox}

For continuous-time signals the sum becomes an integral (this is where the calculus of
Section~\ref{ch42:calculus} enters):
\[
  \inner{x}{y} = \int_{-\infty}^{\infty} x(t)\,\overline{y(t)}\,\dd t.
\]
Same idea, same geometry; we are just adding up infinitely many infinitesimal coordinates.

\begin{figure}[H]
\centering
\begin{tikzpicture}[scale=1.0]
  % axes / basis direction
  \draw[->, gray] (-0.4,0) -- (4.6,0);
  \draw[->, gray] (0,-0.4) -- (0,2.8);
  % unit basis vector e
  \draw[targetgreen, ->, ultra thick] (0,0) -- (2.0,0);
  \node[targetgreen] at (2.0,-0.32) {$\vec e$};
  % signal vector x
  \draw[sourceblue, ->, ultra thick] (0,0) -- (3.2,2.2);
  \node[sourceblue] at (3.35,2.35) {$\vec x$};
  % projection foot
  \draw[dashed, gray] (3.2,2.2) -- (3.2,0);
  % projection length highlighted on the axis
  \draw[specred, ->, very thick] (0,0) -- (3.2,0);
  \node[specred] at (1.6,0.28) {$\inner{\vec x}{\vec e}$};
  % right-angle mark
  \draw (3.0,0) -- (3.0,0.2) -- (3.2,0.2);
\end{tikzpicture}
\caption{The inner product as a projection. With $\vec e$ a unit vector, $\inner{\vec x}{\vec e}$
is the length of the shadow that $\vec x$ casts on the $\vec e$ direction -- how much of $\vec e$
is present in $\vec x$. Every Fourier coefficient is one such shadow, taken onto a sinusoid.}
\label{ch42:fig-projection}
\end{figure}

\subsection{Orthogonality and orthonormal bases}

Two vectors are \emph{orthogonal} when their inner product is zero -- they share nothing, they
are perpendicular. A set of vectors is \emph{orthonormal} when each has length $1$ and every pair
is orthogonal. Orthonormal sets are the gold standard for a coordinate system because they make
decomposition free of bookkeeping. If $\{\vec e_0, \vec e_1, \dots, \vec e_{N-1}\}$ is an
orthonormal basis of $\C^N$, then \emph{any} vector expands as
\[
  \vec x = \sum_{k=0}^{N-1} c_k\,\vec e_k, \qquad c_k = \inner{\vec x}{\vec e_k}.
\]
The coefficient $c_k$ is just the inner product of $\vec x$ with the $k$-th basis vector -- the
projection of $\vec x$ onto that direction. No matrix inversion, no system of equations: project
and read off.

\begin{whybox}
Why does $c_k = \inner{\vec x}{\vec e_k}$? Take the inner product of both sides of $\vec x =
\sum_j c_j \vec e_j$ with $\vec e_k$. By linearity it becomes $\inner{\vec x}{\vec e_k} = \sum_j
c_j \inner{\vec e_j}{\vec e_k}$. Orthonormality means $\inner{\vec e_j}{\vec e_k}$ is $1$ when
$j=k$ and $0$ otherwise, so every term in the sum vanishes except $j=k$, leaving $c_k$. The whole
power of an orthonormal basis is that this collapse happens.
\end{whybox}

\begin{happeningbox}
The Fourier transform \emph{is} a change to an orthonormal basis. The DFT basis vectors are the
complex exponentials $\vec e_k$ with entries $e_k[n] = \frac{1}{\sqrt N}\,e^{\jj 2\pi kn/N}$ -- a
phasor that completes exactly $k$ whole turns over the $N$ samples. They are mutually orthogonal
(a fact you can verify with a geometric-series sum), and the $k$-th DFT coefficient is precisely
$\inner{\vec x}{\vec e_k}$: the amount of frequency $k$ in the signal. ``Analyzing a signal into
frequencies'' and ``writing a vector in an orthonormal basis'' are the same sentence in two
dialects.
\end{happeningbox}

\subsection{Matrices: linear maps on signals}

A matrix is a rule for turning one vector into another that respects addition and scaling:
$\mat A(\vec x + \vec y) = \mat A\vec x + \mat A\vec y$ and $\mat A(c\vec x) = c\,\mat A\vec x$.
Such a rule is called \emph{linear}, and a huge fraction of DSP operations are linear maps:
\begin{itemize}
  \item The DFT itself is multiplication by the $N\times N$ matrix $\mat W$ with entries $W_{kn} =
  e^{-\jj 2\pi kn/N}$. The FFT is just a clever factorization of this one matrix.
  \item An LTI filter, applied to a finite signal, is multiplication by a special matrix (a
  circulant or Toeplitz matrix) built from the impulse response.
  \item Decimation, interpolation, the discrete cosine transform, and orientation rotations of an
  IMU reading are all matrix multiplications.
\end{itemize}
Reading $\vec y = \mat A\vec x$ as ``each output sample is a weighted sum of input samples, with
the weights given by a row of $\mat A$'' is usually the most useful way to see what a matrix does.

\subsection{Eigenvectors: the directions a system leaves alone}

For a square matrix $\mat A$, a nonzero vector $\vec v$ is an \emph{eigenvector} with
\emph{eigenvalue} $\lambda$ if
\[
  \mat A \vec v = \lambda\,\vec v.
\]
In words: the map $\mat A$ does not rotate $\vec v$ into a new direction; it only stretches it by
the factor $\lambda$. Eigenvectors are the ``natural axes'' of a transformation -- the special
inputs that come out parallel to themselves.

\begin{aestheticbox}
Eigenvectors are the deep reason sinusoids matter so much. \textbf{Complex exponentials are the
eigenvectors of every LTI system.} Feed a pure $e^{\jj\omega n}$ into any linear time-invariant
filter and you get back the \emph{same} exponential, merely scaled by a complex number -- the
frequency response $H(e^{\jj\omega})$ at that frequency. The filter cannot change the frequency,
only the amplitude and phase. This is why we analyze signals frequency by frequency: in the
frequency basis, every LTI system becomes diagonal -- it acts independently on each component by
simple multiplication. Convolution in time turns into multiplication in frequency precisely
because the Fourier basis diagonalizes convolution. Eigenvectors are not an abstraction bolted on
later; they are the engine.
\end{aestheticbox}

\begin{example}[The DFT basis is orthogonal]
\label{ch42:ex-dft-orth}
Take $N = 4$ and the two basis phasors $e_1[n] = e^{\jj 2\pi n/4} = \jj^n$ and $e_0[n] = 1$. Then
$\inner{\vec e_1}{\vec e_0} = \sum_{n=0}^{3} \jj^n\cdot\overline 1 = 1 + \jj + \jj^2 + \jj^3 = 1
+ \jj - 1 - \jj = 0$. The four powers of $\jj$ sit at the four corners of the unit square and
cancel in pairs. The same telescoping happens for any two distinct DFT frequencies: the partial
sum of a full set of equally spaced points on the circle is zero. That cancellation is
orthogonality, and orthogonality is why the inverse DFT is just another sum, not a matrix
inversion.
\end{example}

\begin{connectionbox}
Inner products, orthonormal bases, projection, matrices as linear maps, and eigen-decomposition
are developed carefully -- with proofs and pictures -- in \emph{Linear Algebra From the Ground
Up}. The chapters on the spectral theorem and on orthogonal projection are the direct backbone of
Fourier analysis. This appendix gives the operating manual; that book gives the theory.
\end{connectionbox}

\begin{goalbox}
The wrist tracker uses linear algebra at two levels. Low level: every windowed FFT of the sensor
stream is a matrix multiply (or its fast factorization), turning a block of samples into a
spectrum. High level: gesture \emph{classification} typically projects a high-dimensional feature
vector (band powers, statistics across the IMU and EMG channels) down to a few discriminative
directions -- principal component analysis (an eigen-decomposition of the covariance matrix) or a
learned linear layer -- before a decision is made. Knowing that a classifier's first stage is
``project onto a handful of well-chosen directions'' demystifies most of the machine-learning
chapter.
\end{goalbox}

\section{The calculus DSP needs}
\label{ch42:calculus}

Discrete signal processing runs on sums, but the \emph{continuous} theory that justifies it --
the continuous-time Fourier transform, the Laplace transform, the sampling theorem, the analog
filters we approximate -- runs on integrals. You do not need the full machinery of calculus to
follow this book, but you need three things solidly: what an integral really is, the two or three
rules for manipulating the ones we meet, and the strange and wonderful delta function.

\subsection{The integral as a continuous sum}

An integral is what a sum becomes when the steps shrink to zero. Chop the time axis into slices
of width $\Delta t$, sample the signal in each slice, multiply value by width, and add:
\[
  \int_a^b x(t)\,\dd t \;=\; \lim_{\Delta t\to 0}\;\sum_n x(t_n)\,\Delta t.
\]
Geometrically this is the signed area under the curve $x(t)$ between $a$ and $b$. The symbol
$\int$ is an elongated S for ``sum,'' and $\dd t$ is the infinitesimal width that used to be
$\Delta t$. Every integral in this book is, at bottom, ``add up the signal continuously.''

\begin{intuitionbox}
Whenever you see an integral in a transform, read it as a weighted average or a running total.
The Fourier transform $X(f) = \int x(t)\,e^{-\jj 2\pi f t}\,\dd t$ is the continuous version of
the inner-product sum from Section~\ref{ch42:linalg}: it multiplies the signal, point by point,
by a probe sinusoid $e^{-\jj 2\pi f t}$ and adds up the result. A large total means the signal
marches in step with that frequency; a total near zero means it does not. Same projection idea,
infinitely fine.
\end{intuitionbox}

\subsection{The few rules we actually use}

Three properties cover almost every integral manipulation in the main text.

\textbf{Linearity.} Integration adds and scales like any good averaging operation:
\[
  \int \big(a\,x(t) + b\,y(t)\big)\,\dd t = a\int x(t)\,\dd t + b\int y(t)\,\dd t.
\]
This is exactly why the Fourier transform is linear, which in turn is why we can analyze a signal
one frequency component at a time.

\textbf{The exponential integral.} The single antiderivative we lean on hardest is that of the
complex exponential, because $\frac{\dd}{\dd t} e^{\jj\omega t} = \jj\omega\,e^{\jj\omega t}$
implies
\[
  \int e^{\jj\omega t}\,\dd t = \frac{1}{\jj\omega}\,e^{\jj\omega t} + C \quad(\omega\neq 0).
\]
Evaluated over a finite window this gives the $\sinc$ that pervades sampling and windowing; over
the whole line (in the limiting, distributional sense) it gives a delta in frequency. Either way,
the integral of a complex exponential is the workhorse of transform calculus.

\textbf{Integration by parts.} For the occasional derivative-times-function integral,
\[
  \int u\,\dd v = uv - \int v\,\dd u,
\]
which is just the product rule for derivatives, read backwards and integrated. It is how one
proves the differentiation property of the Fourier and Laplace transforms (a derivative in time
becomes multiplication by $\jj\omega$ or $s$ in the transform domain).

\begin{example}[Fourier transform of a rectangular pulse]
\label{ch42:ex-rect}
Let $x(t) = 1$ for $\abs{t}\le \tfrac12$ and $0$ otherwise -- the unit rectangle $\rect(t)$. Its
Fourier transform is a single application of the exponential integral:
\[
  X(f) = \int_{-1/2}^{1/2} e^{-\jj 2\pi f t}\,\dd t
  = \left.\frac{e^{-\jj 2\pi f t}}{-\jj 2\pi f}\right|_{-1/2}^{1/2}
  = \frac{e^{\jj\pi f} - e^{-\jj\pi f}}{\jj 2\pi f}
  = \frac{\sin(\pi f)}{\pi f} = \sinc(f).
\]
The last step used the inverse Euler formula for sine. A pulse in time is a $\sinc$ in frequency:
the most important transform pair in sampling theory falls out of one integral and one Euler
substitution.
\end{example}

\begin{figure}[H]
\centering
\begin{tikzpicture}
\begin{axis}[
    width=12cm, height=5cm,
    xlabel={frequency $f$}, ylabel={$\sinc(f)=\dfrac{\sin\pi f}{\pi f}$},
    xmin=-5, xmax=5, ymin=-0.32, ymax=1.1,
    axis lines=middle, samples=200, domain=-5:5,
    xtick={-4,-3,-2,-1,1,2,3,4},
]
\addplot[sourceblue, thick] {sin(deg(pi*x))/(pi*x)};
\addplot[gray, dashed] coordinates {(-5,0) (5,0)};
\end{axis}
\end{tikzpicture}
\caption{The sinc function, the Fourier transform of the unit rectangular pulse. It crosses zero
at every nonzero integer and decays as $1/f$. This single curve governs spectral leakage,
windowing, and ideal-filter ringing throughout the book.}
\label{ch42:fig-sinc}
\end{figure}

\subsection{The delta function: an infinitely thin spike}

We need an idealized impulse: something that is zero everywhere except at one instant, yet has
total area $1$. No ordinary function can do this, so the \emph{Dirac delta} $\delta(t)$ is not a
function in the usual sense but a \emph{distribution} -- defined entirely by how it behaves inside
an integral. Its one defining property is the \emph{sifting property}:
\[
  \int_{-\infty}^{\infty} x(t)\,\delta(t - t_0)\,\dd t = x(t_0).
\]
Sliding a delta to time $t_0$ and integrating against $x$ ``sifts out'' the single value
$x(t_0)$. Picture $\delta$ as the limit of a tall narrow rectangle of area $1$ as its width goes
to zero and its height to infinity; the area is pinned at $1$ no matter how thin it gets.

\begin{trapbox}
Do not ask for the \emph{value} of $\delta(0)$ -- it is not a number, and chasing it leads
nowhere. The delta only ever means something \emph{under an integral}. Treat $\delta$ as a verb
(``sample here''), not a noun (``a value''). Every legitimate use in this book either integrates
against it or convolves with it.
\end{trapbox}

\begin{happeningbox}
The delta is the bridge between continuous and discrete worlds. Sampling a continuous signal is
modeled as multiplying it by an infinite comb of deltas spaced $T_s$ apart, $\sum_n \delta(t -
nT_s)$: each delta sifts out one sample value. Convolving any signal with a single $\delta(t)$
leaves it unchanged -- $\delta$ is the identity for convolution -- which is why the impulse
response (the output when the input is a delta) completely characterizes an LTI system. The
spike that seems too pathological to be a function turns out to be the most useful test signal in
the field.
\end{happeningbox}

\begin{connectionbox}
Limits, the Riemann integral, the fundamental theorem of calculus, and a careful treatment of
distributions are the subject of \emph{Calculus From the Ground Up}. This appendix takes the
integral as ``a continuous sum'' and the delta as ``the sifting rule''; that book proves both
from the foundations.
\end{connectionbox}

\section{Just enough probability for noise}
\label{ch42:probability}

Every real measurement is corrupted. An IMU axis at rest does not read a clean $0$ or $1g$; it
jitters. An EMG electrode picks up thermal noise, motion artifacts, and mains hum. To reason
about how much of a signal we can trust, we describe the random part with probability. We need
only the basics.

\subsection{Random variables, mean, and variance}

A \emph{random variable} is a quantity whose value is uncertain, described by a
\emph{distribution} that says how likely each value is. Two numbers summarize it for most DSP
purposes. The \emph{mean} (or expectation) is the long-run average value,
\[
  \mu = \E[X],
\]
and the \emph{variance} is the average squared deviation from the mean, a measure of spread,
\[
  \sigma^2 = \Var[X] = \E\big[(X - \mu)^2\big] = \E[X^2] - \mu^2.
\]
Its square root $\sigma$, the \emph{standard deviation}, has the same units as the signal and is
the natural ``size'' of the noise. For a noise process we usually want $\mu = 0$ (no constant
bias) and a small $\sigma$ (little jitter).

\begin{connectionbox}
Notice that the variance is an inner product in disguise: it is the expected value of a product,
$\E[(X-\mu)^2]$, exactly as energy was $\inner{\vec x}{\vec x}$ in Section~\ref{ch42:linalg}.
This is not a coincidence. \emph{Probability From the Ground Up} develops expectation as an
integral against a density, $\E[g(X)] = \int g(x)\,p(x)\,\dd x$, which makes the whole subject a
special case of the calculus and inner-product machinery you already have. Mean is a projection;
variance is a squared length.
\end{connectionbox}

\subsection{The Gaussian, and white noise}

The single most important distribution for noise is the \emph{Gaussian} (normal) distribution,
with density
\[
  p(x) = \frac{1}{\sqrt{2\pi}\,\sigma}\,\exp\!\left(-\frac{(x-\mu)^2}{2\sigma^2}\right).
\]
Its bell shape is set entirely by $\mu$ (where it is centered) and $\sigma$ (how wide). It earns
its dominance through the \emph{central limit theorem}: when many small independent disturbances
add up -- as they do in a real sensor -- their sum is approximately Gaussian regardless of the
details. That is why thermal and electronic noise is so often modeled as Gaussian.

\begin{figure}[H]
\centering
\begin{tikzpicture}
\begin{axis}[
    width=11cm, height=5cm,
    xlabel={value $x$}, ylabel={density $p(x)$},
    xmin=-4, xmax=4, ymin=0, ymax=0.85,
    axis lines=left, samples=160, domain=-4:4,
    legend pos=north east, legend cell align=left,
]
\addplot[sourceblue, thick] {1/(sqrt(2*pi))*exp(-x^2/2)};
\addlegendentry{$\sigma=1$}
\addplot[specred, thick, dashed] {1/(0.6*sqrt(2*pi))*exp(-x^2/(2*0.36))};
\addlegendentry{$\sigma=0.6$}
\end{axis}
\end{tikzpicture}
\caption{Gaussian densities with mean $0$. Smaller $\sigma$ means less spread (a sharper, taller
bell) -- less noise. About $68\%$ of samples fall within $\pm\sigma$ of the mean and $95\%$
within $\pm 2\sigma$.}
\label{ch42:fig-gaussian}
\end{figure}

\emph{White noise} is a sequence of random samples that are uncorrelated from one to the next and
share the same distribution -- typically zero-mean Gaussian. ``White'' refers to its flat power
spectrum: like white light, it contains every frequency in equal measure. It is the standard
model for the irreducible jitter floor of a sensor, and the baseline against which we measure how
well a filter cleans a signal.

\subsection{Signal-to-noise ratio}

The one figure of merit that ties probability back to engineering is the \emph{signal-to-noise
ratio}, the ratio of signal power to noise power, almost always quoted in decibels:
\[
  \mathrm{SNR} = \frac{P_{\text{signal}}}{P_{\text{noise}}}, \qquad
  \mathrm{SNR}_{\text{dB}} = 10\log_{10}\!\frac{P_{\text{signal}}}{P_{\text{noise}}}.
\]
Noise power is just the variance $\sigma^2$ of a zero-mean noise process. A higher SNR means a
cleaner signal; averaging $N$ independent measurements of a static quantity reduces noise power
by a factor of $N$ (improving SNR by $10\log_{10} N$ dB), which is the quantitative reason
averaging works.

\begin{example}[Averaging beats down noise]
\label{ch42:ex-averaging}
Suppose a stationary forearm gives an accelerometer reading $x[n] = g + w[n]$, where $g$ is the
true gravity component and $w[n]$ is zero-mean white noise with standard deviation $\sigma =
0.02g$. A single sample has noise power $\sigma^2$. Average $N = 100$ samples: the signal $g$ is
unchanged, but the noise variance drops to $\sigma^2/100$, so the standard deviation falls by
$\sqrt{100} = 10$, to $0.002g$. The SNR improves by $10\log_{10}100 = 20$ dB. This is exactly the
trick a stillness detector or a slow-tilt estimator uses on the wrist device -- and it costs only
a small averaging window.
\end{example}

\begin{connectionbox}
Distributions, expectation as an integral, independence, the central limit theorem, and random
processes are built from the ground up in \emph{Probability From the Ground Up}. The noise and
spectral-estimation chapters of this book (Chapters~20--21) use those ideas in earnest; this
section is only the vocabulary needed to start.
\end{connectionbox}

\begin{goalbox}
A clear-eyed view of the wrist tracker starts with noise. The reason a wrist-worn IMU \emph{alone}
cannot recover individual finger joint angles is partly geometric -- the fingers move
distally, far from the wrist sensor -- and partly statistical: the tiny mechanical couplings that
do reach the wrist are buried below the sensor's noise floor, and no amount of averaging recovers
a signal that was never above $\sigma$ in the first place. Coarse, discrete wrist and
hand gestures \emph{are} detectable because they produce large, repeatable signals well above the
noise. Continuous finger pose needs sensors closer to the source -- EMG or force myography over
the forearm muscles, or rings on the fingers themselves. Probability is what lets you state that
limit precisely instead of hoping: if the desired signal is below the noise floor, it is gone.
See Chapters~34 and~37 for the sensor-by-sensor accounting.
\end{goalbox}

\section{Exercises}
\label{ch42:exercises}

\begin{exercise}
Write $z = -1 + \jj\sqrt 3$ in polar form $r e^{\jj\theta}$ with $\theta\in(-\pi,\pi]$. Then
compute $z^3$ in polar form and convert the result back to rectangular form. (Hint:
$\operatorname{atan2}$ and the quadrant matter.)
\end{exercise}

\begin{exercise}
Using the inverse Euler formulas, prove the identity $\cos^2\theta = \tfrac12 + \tfrac12\cos
2\theta$ by writing $\cos\theta = (e^{\jj\theta} + e^{-\jj\theta})/2$, squaring, and simplifying.
Explain in one sentence what this identity says about the spectrum of a squared sinusoid.
\end{exercise}

\begin{exercise}
Let $\vec x = (1, \jj, -1, -\jj)$ in $\C^4$. Compute $\norm{\vec x}^2$ using the complex inner
product, being careful to conjugate. Then compute $\inner{\vec x}{\vec e_1}$ where $e_1[n] =
e^{\jj 2\pi n/4} = \jj^n$. What does the result tell you about the frequency content of $\vec x$?
\end{exercise}

\begin{exercise}
Show directly that the two DFT-basis vectors of length $N=4$, $e_1[n] = \jj^n$ and $e_3[n] =
\jj^{3n} = (-\jj)^n$, are orthogonal under the complex inner product. Why is conjugation
essential to getting zero here?
\end{exercise}

\begin{exercise}
A complex exponential $\vec v$ with $v[n] = e^{\jj\omega n}$ is claimed to be an eigenvector of
any LTI filter. Take the simplest filter, a one-sample delay $y[n] = x[n-1]$. Verify that feeding
$v[n]$ in produces $\lambda\,v[n]$ and find the eigenvalue $\lambda$ in terms of $\omega$. What
are $\abs{\lambda}$ and $\arg(\lambda)$, and what do they mean physically?
\end{exercise}

\begin{exercise}
Compute $\int_{-T/2}^{T/2} e^{-\jj 2\pi f t}\,\dd t$ for a window of width $T$ and show it equals
$T\,\sinc(fT)$. At which frequencies $f$ does it vanish? Relate this to why a finite observation
window smears a single tone across nearby frequencies (spectral leakage).
\end{exercise}

\begin{exercise}
Use the sifting property to evaluate $\int_{-\infty}^{\infty}(t^2 + 3)\,\delta(t-2)\,\dd t$ and
$\int_{-\infty}^{\infty}\cos(\pi t)\,\delta(t - \tfrac12)\,\dd t$. Then explain, in one or two
sentences, why convolving an arbitrary signal $x(t)$ with $\delta(t - t_0)$ shifts it to
$x(t - t_0)$.
\end{exercise}

\begin{exercise}
A zero-mean white-noise sensor channel has standard deviation $\sigma = 0.05g$. (a) What is its
noise power? (b) If you average $N$ consecutive samples of a stationary reading, by what factor
does the noise standard deviation drop? (c) How many samples must you average to bring the noise
standard deviation down to $0.01g$, and by how many decibels does that improve the SNR?
\end{exercise}

\begin{exercise}
For a Gaussian random variable with mean $\mu$ and variance $\sigma^2$, roughly what fraction of
samples fall within $\pm 2\sigma$ of the mean? If a stillness detector flags ``moving'' whenever
a sample exceeds $\pm 2\sigma$ of the resting distribution, what false-alarm rate would you
expect when the wrist is in fact still, and is that acceptable for triggering a MIDI note?
\end{exercise}

\begin{exercise}
(Synthesis.) Explain, in a short paragraph each, how each of the four topics in this appendix --
complex numbers, inner products/orthonormal bases, integration/the delta, and probability --
appears in the single act of computing one windowed FFT of an IMU stream and reading off the
dominant motion frequency for gesture detection. Be specific about where each piece enters.
\end{exercise}
